\section{Treat a pointer, a capacity and committed output as separate facts}
The backing word and byte arenas are larger than many calls. Source length authorizes
reads; destination capacity bounds potential output; the final committed prefix bounds
permitted stores. None of these quantities is interchangeable. A mapped address outside
the relevant interval still violates the interface.

Observe writes during execution, rather than comparing final memory alone. A failed
candidate must not leave a partial destination write in the history. The caller's
source and metadata also remain read-only until the function returns. The supplied
caller then stores the two counts and status, so inspect the function at
\texttt{codec\_done} before those caller stores. A direct grader call may instead
supply \texttt{function\_return}; return through \texttt{ra} in either case.

Both byte-source and byte-destination pointers may be unaligned. Byte loads and stores
make their boundaries explicit. A signed byte load can also work if every use masks
the payload and continuation bits appropriately; an unsigned load is a straightforward
choice. Numerical grading should accept such equivalent operand handling rather than
requiring an instruction spelling.

Each call starts with predecessor zero, even in one continuing simulator session.
Encode, decode, run a zero-capacity failure, and then encode a new sequence without
restarting the simulator. No previous reading or partial unsigned accumulator should
survive as an implicit input. Saved registers and the full stack pointer must return
to their incoming values on success and every error path.
